\section{Nonvanishing and the functional equations}
\label{sec:deduction}

We return to the character with which the introduction began. Proposition~\ref{prop:arithmetic}, proved using Theorem~\ref{hyp:uniform-arithmetic-data}, supplies the intervals required by the Laplace-transform criterion for each primitive character. We pass to arbitrary moduli using the finite
Euler factors, then use the functional equation to relate \(s\) to \(1-s\)
in the open critical strip. For \(\zeta\) the
explicit functional equation also determines which zeros can
occur to the left of the strip. These steps prove all parts of
Theorem~\ref{thm:claims}.

\begin{definition}[Completion at the primitive conductor]\label{def:completed-dirichlet}
For a primitive Dirichlet character \(\chi\) of conductor \(f_\chi\),
let \(a_\chi=(1-\chi(-1))/2\). Thus \(a_\chi=0\) for an even character
and \(a_\chi=1\) for an odd character. Its completed \(L\)-function is
\[
 \Lambda(s,\chi)=
 \left(\frac{f_\chi}{\pi}\right)^{(s+a_\chi)/2}
 \Gamma\!\left(\frac{s+a_\chi}{2}\right)L_{f_\chi}(s,\chi).
\]
In this chapter \(q\) is the modulus of the character under consideration;
for an induced character the completion above is applied to its primitive
inducer \(\chi^*\). The omitted local factors are those of
\eqref{eq:imprimitive-euler-factors}.
\end{definition}

\begin{construction}
\begin{tabular}{@{}l@{\hspace{1.2em}}l@{}}
\toprule
\textit{Characters and domain} & \textit{Mathematical input} \\
\midrule
Primitive \(\psi,\ \RePart s>1/2\) & Weighted sum on \([Y,Y^{3/2}]\) \\
Induced \(\chi,\ \RePart s>0\) & Extra Euler factors are nonzero \\
Primitive \(\chi,\ 0<\RePart s<1\) & \((s,\chi)\mapsto(1-s,\bar\chi)\) \\
\(\zeta,\ \RePart s<1/2,\ s\ne0\) & Cosine factor, with \(\zeta(1-s)\ne0\) \\
\bottomrule
\end{tabular}
\caption{Domains of the final deductions. The first row excludes \(s=1\) for the principal character.}\label{fig:analytic-scopes}
\end{construction}

\begin{proof}[Proof of Theorem~\ref{thm:claims}]
\proofstep{The right half-plane.}
Fix a primitive character
\(\psi\) and a proposed zero \(\rho\) with \(\RePart\rho>1/2\).
Choose \(1/2<\theta<\min(1,\RePart\rho)\), and then take
\(0<\epsilon<\epsilon_0\) with \(9\epsilon\leq\theta-1/2\).
On the intervals supplied by Proposition~\ref{prop:arithmetic}, with \(Y\) arbitrarily large,
\[
   G_{\mathcal K}(y)\le y^{1/2+9\epsilon}\le y^\theta
            \qquad(Y\le y\le Y^{3/2}).
\]
The weights of \(\mathcal K\) are positive and \(\psi\in\mathcal K\),
so Theorem~\ref{thm:window} with \(\delta=1/2\) excludes \(\rho\).

This proves the primitive case at the conductor. For a character \(\chi\bmod q\)
induced by \(\chi^*\) of conductor \(f_\chi\mid q\), the Euler-factor
identity \cite[(25.15.4)]{dirichlet} is
\[
 L_q(s,\chi)=L_{f_\chi}(s,\chi^*)\prod_{\substack{p\mid q\\p\nmid f_\chi}}
       (1-\chi^*(p)p^{-s})
\]
by comparison of the Euler products and meromorphic continuation.
The factors are nonzero for \(\RePart s>0\) since
\(\lvert\chi^*(p)p^{-s}\rvert\leq p^{-\RePart s}<1\). This gives the
nonvanishing for \(\RePart s>1/2\), with \(s=1\) excluded for the principal character.

\proofstep{The open critical strip.}
We can now reflect the nonvanishing statement inside the strip. For a primitive
nonprincipal \(\chi\), let \(\bar\chi\) denote its
complex conjugate and let \(\varepsilon(\chi)\) denote its root number. With the completion in
Definition~\ref{def:completed-dirichlet}, the primitive Dirichlet
functional equation \cite[(25.15.5)]{dirichlet} takes the form
\[
\Lambda(s,\chi)=\varepsilon(\chi)\Lambda(1-s,\bar\chi),
                   \qquad |\varepsilon(\chi)|=1.
\]
For \(\chi=1\) we have \(f_\chi=1\) and \(a_\chi=0\). The functional
equation of the Riemann \(\xi\)-function gives
\[
 \Lambda(s,1)=\pi^{-s/2}\Gamma(s/2)\zeta(s)=\Lambda(1-s,1)
                     \qquad (0<\RePart s<1),
\]
after cancelling its factor \(s(s-1)/2\).
\cite[(25.4.3)--(25.4.4)]{zeta}

In both cases the gamma factors are finite and nonzero there.
Thus for primitive characters the functional equations give
\[
 L(s,\chi)=0\quad\Longleftrightarrow\quad
 L(1-s,\bar\chi)=0\qquad(0<\RePart s<1).
\]
If \(\RePart s<1/2\), then \(\RePart(1-s)>1/2\), which
contradicts primitive nonvanishing for \(\bar\chi\).
Primitive nonvanishing also excludes \(\RePart s>1/2\).
For an induced character, the Euler factors above are units
at every point of the open strip, so its zeros there agree with the
zeros of the primitive character inducing it.

\proofstep{The remaining zeta zeros.}
For \(\RePart s<1/2\) and \(s\ne 0\), the zeta functional
equation \cite[(25.4.1)]{zeta} gives
\[
 \zeta(s)=2(2\pi)^{s-1}\Gamma(1-s)
           \cos\left(\frac{\pi(1-s)}2\right)\zeta(1-s).
\]
Here \(\Gamma(1-s)\) is finite and nonzero and the
right-half-plane result gives \(\zeta(1-s)\ne0\).
The point \(s=0\) was excluded, so \(\zeta(1-s)\) is also finite.
Consequently for \(\RePart s<1/2\) and \(s\ne0\),
\[
 \zeta(s)=0 \quad\Longleftrightarrow\quad
 \cos\!\left(\frac{\pi(1-s)}2\right)=0
 \quad\Longleftrightarrow\quad s\in\{-2,-4,-6,\ldots\}.
\]
Together with \(\zeta(0)=-1/2\), this identifies all zeros
strictly to the left of the line. The modulus-one
right-half-plane result excludes the other side away from \(s=1\).
\end{proof}

